\section{Technical Overview}
\label{sec:overview}

\subsection{MLE Oracles for Product Codes from Simpler MLE Oracles}


\subsection{Structured Code Switching for Interleaved Code}
Let's state all in 

\clearpage
\begin{figure}[!t]
\centering
\resizebox{0.75\linewidth}{!}{
\begin{tikzpicture}[x=1cm,y=1cm]

% Colors (approximate to the figure)
\definecolor{dashblue}{RGB}{34,83,167}
\definecolor{dashgreen}{RGB}{30,166,103}
\definecolor{cream}{RGB}{247,239,224}
\definecolor{mint}{RGB}{183,232,210}
\definecolor{sky}{RGB}{212,233,246}
\definecolor{stroke}{RGB}{70,70,70}


% --- Colors (same style as yours) ---
\definecolor{bgblue}{RGB}{235, 245, 255}
\definecolor{bggreen}{HTML}{F2FFCE}
\definecolor{nodeblue}{RGB}{212,233,246}
\definecolor{nodecream}{RGB}{247,236,214}
\definecolor{nodegreen}{RGB}{221,244,234}
\definecolor{borderblue}{RGB}{100, 149, 237}
\definecolor{bordergreen}{RGB}{46, 125, 50}
\definecolor{textblue}{RGB}{0, 51, 102}
\definecolor{arrowgray}{RGB}{128, 128, 128}
\definecolor{edgegray}{RGB}{140,140,140}

\definecolor{rainbowyellow}{HTML}{FFF3CE}
\definecolor{rainbowgreen}{HTML}{D6E9D5}
\definecolor{rainbowred}{HTML}{FAD8D4}
\definecolor{rainbowcyan}{HTML}{AFE3E6}
\definecolor{rainbowblue}{HTML}{D9E8FC}
\definecolor{rainbowgray}{HTML}{BAC8D3}
\definecolor{rainbowpurple}{HTML}{D5D6E9}
\definecolor{rainbowviolet}{HTML}{E6AFE3}
\definecolor{rainboworange}{HTML}{FFDACE}
\definecolor{rainbowsky}{HTML}{99CCEF}

\definecolor{titleblue}{HTML}{192841}

% Parameters
\def\dx{0.92}
\def\dy{0.90}
\def\r{0.3}
\def\left{0}
\def\top{0}
\pgfmathsetmacro\rows{7.0}
\pgfmathsetmacro\cols{7.0}
\pgfmathsetmacro{\ybottom}{-(\rows-1)*\dy - 0.70}


% Blaze illustration
\begin{scope}
% Helper: circle fill by row/col
% Row index i = 1..7 from top to bottom; col j = 1..7 from left to right.
\foreach \i in {1,...,7} {
  % y coordinate for row
  \pgfmathsetmacro{\yy}{-(\i-1)*\dy}

  % Long dashed blue rounded rectangle for each row
   \pgfmathtruncatemacro{\drawenc}{(\i==5) || (\i==7)}
   \ifnum \drawenc=1
     \draw[draw=dashblue, line width=1.1pt, dash pattern=on 5pt off 2pt on 1pt off 2pt, fill=rainbowcyan,
        rounded corners=8pt]
    (-0.45, \yy+0.38) rectangle (3.2, \yy-0.38);
    \fi

    \node[anchor=west, font=\large\selectfont] at (-2.3, \yy)         {$\Root_{\i} \Leftarrow \bm{c}_{\i}^{\,\trans} $};

    \draw[dashed] (-5,\yy-0.45) -- (3.5,\yy-0.45);

  % Circles
  \foreach \j in {1,...,4} {
        \pgfmathsetmacro{\xx}{(\j-1)*\dx}
    
        % Decide fill color
        % First 3 columns: top 3 rows cream, bottom 4 rows mint.
        % Last 4 columns: all sky.
          \ifnum\i<5
            \def\fillcol{cream}
          \else
            \def\fillcol{mint}
            \node[anchor=center, font=\large\selectfont] at (\xx, \yy) {$c_{\i\j}$};
          \fi

          
        \draw[draw=stroke, line width=1pt, fill=\fillcol] (\xx, \yy) circle (\r);

        \ifnum\i<5
            \node[anchor=center, font=\large\selectfont] at (\xx, \yy) {$c_{\i\j}$};
        \else
            \node[anchor=center, font=\large\selectfont] at (\xx, \yy) {$c_{\i\j}$};
        \fi

        \ifnum\j=1
       \node[anchor=east, font=\large\selectfont] at (-3.0, \yy) {$=\langle \bm{c}_{\i}^{\,\trans}, \tilde{\bm{r}} \rangle$};
       \fi     
    }

  % Circles
  \foreach \j in {-4.8} {
        \pgfmathsetmacro{\xx}{(\j-1)*\dx}
    
        % Decide fill color
        % First 3 columns: top 3 rows cream, bottom 4 rows mint.
        % Last 4 columns: all sky.
          \ifnum\i<5
            \def\fillcol{mint}
          \else
            \ifnum\i=5
            \def\fillcol{green}
            \else
            \ifnum\i=7
            \def\fillcol{green}
            \else
            \def\fillcol{mint}
            \fi\fi
          \fi
          
        \draw[draw=stroke, line width=1pt, fill=\fillcol] (\xx, \yy) circle (\r);
        \node[anchor=center, font=\large\selectfont] at (\xx, \yy) {$u_{\i}$};
        \node[anchor=center, font=\large\selectfont] at (\xx, 1.5) {};
        \node[anchor=center, font=\large\selectfont] at (\xx, 0.8) {};
    }
}



% Three vertical dashed green rounded rectangles over the first 3 columns
\foreach \j in {1,...,4} {
  \pgfmathsetmacro{\xleft}{(\j-1)*\dx - 0.36}
  \pgfmathsetmacro{\xright}{(\j-1)*\dx + 0.36}

  \draw[
    draw=dashgreen,
    line width=1.1pt,
    dash pattern=on 5pt off 2pt on 1pt off 2pt,
    rounded corners=8pt
  ]
    (\xleft,0.45) rectangle (\xright,\ybottom);
}

% Three vertical dashed green rounded rectangles over the first 3 columns
\foreach \j in {-4.8} {
  \pgfmathsetmacro{\xleft}{(\j-1)*\dx - 0.36}
  \pgfmathsetmacro{\xright}{(\j-1)*\dx + 0.36}

  \draw[
    draw=dashgreen,
    line width=1.1pt,
    dash pattern=on 5pt off 2pt on 1pt off 2pt,
    rounded corners=8pt
  ]

  
    (\xleft,0.45) rectangle (\xright,\ybottom);
}

% ============
% ============ Draw the left basefold code
\foreach \i in {1,...,8} {
    % y coordinate for row
  \pgfmathsetmacro{\yy}{-(\i-1)*\dy}
  % Circles
  \foreach \j in {-6.0} {
        \pgfmathsetmacro{\xx}{(\j-1)*\dx}
            
            \ifnum\i=1
            \draw[draw=dashblue, line width=1.1pt, dash pattern=on 5pt off 2pt on 1pt off 2pt, fill=rainbowgray,
        rounded corners=8pt]
            (\xx+0.38, -6.95 ) rectangle (\xx-0.38, 0.5);
            \fi
            
        % Decide fill color
        % First 3 columns: top 3 rows cream, bottom 4 rows mint.
        % Last 4 columns: all sky.
          \ifnum\i<5
            \def\fillcol{rainbowred}
          \else
            \def\fillcol{rainbowsky}
          \fi

        \draw[draw=stroke, line width=1pt, fill=\fillcol] (\xx, \yy) circle (\r);
        \node[anchor=center, font=\large\selectfont] at (\xx, \yy) {$c_{\i}'$};
    }
}


% Bottom text: Enc^{\otimes \epsilon^{-1}} \ldots Enc^{\otimes \epsilon^{-1}}
\node[anchor=center, font=\large\selectfont] at (1.4, 1.5) {$\Root\Leftarrow [\Root_1,\ldots,\Root_7]$};
\node[anchor=center, font=\large\selectfont] at (-6.5, 1.5) {$\Root' \Leftarrow \bm{c}' $};
\node[anchor=center, font=\large\selectfont] at (1.5, -6.55) {$\Enc~~\Enc~~\Enc~~\Enc$};
\node[anchor=center, font=\large\selectfont] at (-4.9, -6.55) {$\bm{u} = \Enc(\bm{v})$};
\node[anchor=center, font=\large\selectfont] at (-5.1, -7.05) {$v = \langle \bm{v}, \bm{l} \rangle$};
\node[anchor=center, font=\large\selectfont] at (-6.4, 0.75) {$\bm{c}' = \Enc'(\bm{v})$};
\node[anchor=center, font=\large\selectfont] at (1.5, 0.9) {$t$};
\node[anchor=center, font=\large\selectfont] at (3.6, -0.9) {$k$};
\node[anchor=center, font=\large\selectfont] at (-7.5, -1.2) {$\bm{v}$};
\node[anchor=center, font=\large\selectfont] at (-2.65, -2.7) {$\mu$};

\node[anchor=center, font=\large\selectfont] at (5.7, 1.4) {$v = \bm{l}^{\,\trans}\bm{M}\tilde{\tilde{\bm{r}}}$};
\node[anchor=center, font=\large\selectfont] at (5.7, 0.6) {$\bm{C} = \bm{G}\bm{M}$};
\node[anchor=center, font=\large\selectfont] at (5.25, -0.2) {$\Root = \Merkle(\bm{C})$};
\node[anchor=center, font=\large\selectfont] at (6.5, -0.8) {$= \Merkle([\Root_1, \ldots, \Root_7])$};

\node[anchor=center, font=\Large\selectfont\bfseries, color=titleblue] at (-1.5, 2.0) {Blaze+};

\draw[-Stealth] (-1,0.7) -- (-4,0.7)
  node[midway, above, font=\Large] {$\bm{v} = \bm{M}\tilde{\bm{r}}$};
\draw[-Stealth] (-1,-7.0) -- (-4,-7.0)
  node[midway, above, font=\Large] {$\bm{u} = \bm{C}\tilde{\bm{r}}$};


%%% Draw braces
\draw[
  decorate,
  decoration={brace, amplitude=6pt}
] (-0.3,0.5) -- (3.1,0.5);
\draw[
  decorate,
  decoration={brace, amplitude=6pt}
] (3.2,0.4) -- (3.2,-3.0);
\draw[
  decorate,
  decoration={brace, amplitude=6pt}
] (-2.2,-5.9) -- (-2.2,0.4);
\draw[
  decorate,
  decoration={brace, amplitude=6pt}
] (-7.0,-2.9) -- (-7.0,0.4);
\end{scope}

%%%%%%%%%%%%%%%%%%%%%%%%%%%%%%%%%%%%%%%%%%%%%%%%%%%%%%%%%%%%%%%%%%%%

% 扁扁的下三角形
\draw[fill=gray!30, draw=black]
(-7,-8) -- (-6.5,-7.5) -- (5,-8);


%%%%%%%%%%%%%%%%%%%%%%%%%%%%%%%%%%%%%%%%%%%%%%%%%%%%%%%%%%%%%%%%%%%%
%%%=================================================================
% DualFold illustration
\begin{scope}[shift={(0.0,-10.5)}]
\pgfmathsetmacro\rows{6.0}
\pgfmathsetmacro\cols{6.0}
\pgfmathsetmacro{\ybottom}{-(\rows-1)*\dy - 0.70}
% Helper: circle fill by row/col
% Row index i = 1..7 from top to bottom; col j = 1..7 from left to right.
%%%%%%%%%%%%%%%
\foreach \i in {1,...,7} {
  % y coordinate for row
  \pgfmathsetmacro{\yy}{-(\i-1)*\dy}

  \draw[draw=dashblue, line width=1.1pt, dash pattern=on 5pt off 2pt on 1pt off 2pt, rounded corners=8pt]
    (-0.45, \yy+0.38) rectangle (5.1, \yy-0.38);

 

  % Long dashed blue rounded rectangle for each row
   \pgfmathtruncatemacro{\drawenc}{(\i==3) || (\i==5)}
   \ifnum \drawenc=1
     \draw[draw=dashblue, line width=1.1pt, dash pattern=on 5pt off 2pt on 1pt off 2pt, fill=rainbowcyan,
        rounded corners=8pt]
    (-0.45, \yy+0.38) rectangle (5.1, \yy-0.38);
    \fi
}

\foreach \j in {2,4} {
  \pgfmathsetmacro{\xleft}{(\j-1)*\dx - 0.36}
  \pgfmathsetmacro{\xright}{(\j-1)*\dx + 0.36}
      \draw[
        draw=dashgreen,
        line width=1.1pt,
        dash pattern=on 5pt off 2pt on 1pt off 2pt,
        rounded corners=8pt,
        % fill=blue!10!rainbowcyan
      ]
        (\xleft,0.45) rectangle (\xright,\ybottom);
}

\foreach \i in {1,...,7} {
  % y coordinate for row
  \pgfmathsetmacro{\yy}{-(\i-1)*\dy}

  \draw[draw=dashblue, line width=1.1pt, dash pattern=on 5pt off 2pt on 1pt off 2pt, rounded corners=8pt]
    (-0.45, \yy+0.38) rectangle (5.1, \yy-0.38);

 

  % Long dashed blue rounded rectangle for each row
   \pgfmathtruncatemacro{\drawenc}{(\i==3) || (\i==5)}
   \ifnum \drawenc=1
     \draw[draw=dashblue, line width=1.1pt, dash pattern=on 5pt off 2pt on 1pt off 2pt,
        rounded corners=8pt]
    (-0.45, \yy+0.38) rectangle (5.1, \yy-0.38);
    \fi

    %%%%%%%%%%% Draw blue circles
       \foreach \j in {3,...,6} {
        \pgfmathsetmacro{\xx}{(\j-1)*\dx}
        \def\fillcol{rainbowblue}
        \draw[draw=stroke, line width=1pt, fill=\fillcol] (\xx, \yy) circle (\r);
          \node[anchor=center, font=\large\selectfont] at (\xx, \yy) {$c_{\i\j}$};
      }

    
      \node[anchor=east, font=\large\selectfont] at (-3.0, \yy) {$=\langle \bm{c}_{\i}^{\,\trans}, \tilde{\bm{r}} \rangle$};
      \node[anchor=east, font=\large\selectfont] at (6.0, \yy) {$\Enc$};
    %%%%%%%%%%%%%%

    \node[anchor=west, font=\large\selectfont] at (-2.3, \yy)         {$\Root_{\i} \Leftarrow \bm{c}_{\i}^{\,\trans} $};

    \draw[dashed] (-5,\yy-0.45) -- (3.5,\yy-0.45);


  % Circles
  \foreach \j in {1,...,2} {
        \pgfmathsetmacro{\xx}{(\j-1)*\dx}
    
        % Decide fill color
        % First 3 columns: top 3 rows cream, bottom 4 rows mint.
        % Last 4 columns: all sky.
          \ifnum\i<4
            \def\fillcol{rainbowred}
          \else
            \def\fillcol{mint}
            \node[anchor=center, font=\large\selectfont] at (\xx, \yy) {$c_{\i\j}$};
          \fi

          
        \draw[draw=stroke, line width=1pt, fill=\fillcol] (\xx, \yy) circle (\r);

        \ifnum\i<4
            \node[anchor=center, font=\large\selectfont] at (\xx, \yy) {$c_{\i\j}$};
        \else
            \node[anchor=center, font=\large\selectfont] at (\xx, \yy) {$c_{\i\j}$};
        \fi
            
    }

    %%%%%%%%%%% Draw selected blue
    \foreach \j in {2,4}{

        
        \pgfmathsetmacro{\drawpurple}{ ((\i==3) && (\j==4)) || ( (\i==5) && (\j==2))}
        % \pgfmathsetmacro{\drawpurple}{ (\i==3) || (\i==5) }
        \ifnum\drawpurple=1
            \pgfmathsetmacro{\xx}{(\j-1)*\dx}
            \def\fillcol{blue!20!white}
            \draw[draw=stroke, line width=1pt, fill=\fillcol] (\xx, \yy) circle (\r);
              \node[anchor=center, font=\large\selectfont] at (\xx, \yy) {$c_{\i\j}$};
          \fi
      }

  % Circles
  \foreach \j in {-4.8} {
        \pgfmathsetmacro{\xx}{(\j-1)*\dx}
    
        % Decide fill color
        % First 3 columns: top 3 rows cream, bottom 4 rows mint.
        % Last 4 columns: all sky.
          \ifnum\i=3
            \def\fillcol{green}
          \else
            \ifnum\i<4
            \def\fillcol{mint}
            \else
            \ifnum\i=5
            \def\fillcol{green}
            \else
            \def\fillcol{mint}
            \fi\fi
          \fi
          
        \draw[draw=stroke, line width=1pt, fill=\fillcol] (\xx, \yy) circle (\r);
        \node[anchor=center, font=\large\selectfont] at (\xx, \yy) {$u_{\i}$};
        \node[anchor=center, font=\large\selectfont] at (\xx, 1.5) {};
        \node[anchor=center, font=\large\selectfont] at (\xx, 0.8) {};
    }
}



% Three vertical dashed green rounded rectangles over the first 6 columns
\foreach \j in {1,...,6} {
  \pgfmathsetmacro{\xleft}{(\j-1)*\dx - 0.36}
  \pgfmathsetmacro{\xright}{(\j-1)*\dx + 0.36}

  \draw[
    draw=dashgreen,
    line width=1.1pt,
    dash pattern=on 5pt off 2pt on 1pt off 2pt,
    rounded corners=8pt
  ]
    (\xleft,0.45) rectangle (\xright,\ybottom - 0.9);
}



% Linear Combination
\foreach \j in {-4.8} {
  \pgfmathsetmacro{\xleft}{(\j-1)*\dx - 0.36}
  \pgfmathsetmacro{\xright}{(\j-1)*\dx + 0.36}

  \draw[
    draw=dashgreen,
    line width=1.1pt,
    dash pattern=on 5pt off 2pt on 1pt off 2pt,
    rounded corners=8pt
  ]

  
    (\xleft,0.45) rectangle (\xright,\ybottom - 0.9);
}


% ============
% ============ Draw the left basefold code
\foreach \i in {1,...,10} {
    % y coordinate for row
  \pgfmathsetmacro{\yy}{-(\i-1)*\dy}
  % Circles
  \foreach \j in {-6.0} {
        \pgfmathsetmacro{\xx}{(\j-1)*\dx}
            
            \ifnum\i=1
            \draw[draw=dashblue, line width=1.1pt, dash pattern=on 5pt off 2pt on 1pt off 2pt, fill=rainbowgray,
            rounded corners=8pt]
            (\xx+0.38, -8.65 ) rectangle (\xx-0.38, 0.5);
            \fi
            
        % Decide fill color
        % First 3 columns: top 3 rows cream, bottom 4 rows mint.
        % Last 4 columns: all sky.
          \ifnum\i<4
            \def\fillcol{rainbowviolet}
          \else
            \def\fillcol{rainbowsky}
          \fi

        \draw[draw=stroke, line width=1pt, fill=\fillcol] (\xx, \yy) circle (\r);
        \node[anchor=center, font=\large\selectfont] at (\xx, \yy) {$c_{\i}'$};
    }
}


\begin{scope}[yshift=-10pt]
%Draw the bottom linear combination
% Long dashed blue rounded rectangle for each row
\foreach \i in {-7,-8}{
    \pgfmathsetmacro{\yy}{(\i-1)*\dy}
     \draw[draw=dashblue, line width=1.1pt, dash pattern=on 5pt off 2pt on 1pt off 2pt, fill=rainbowcyan,
        rounded corners=8pt]
    (-0.45, \yy+0.38) rectangle (5.1, \yy-0.38);

    \ifnum\i=-7
    \node[anchor=west, font=\large\selectfont] at (-2.4, \yy) {$\Root_3 \Leftarrow \bm{c}_3^{\,\trans} $};
    \else
    \node[anchor=west, font=\large\selectfont] at (-2.4, \yy) {$\Root_5 \Leftarrow \bm{c}_5^{\,\trans} $};
    \fi

     \node[anchor=east, font=\large\selectfont] at (6.0, \yy) {$\Enc$};
     
    %%%%%%%%%%% Draw blue circles
       \foreach \j in {1,...,6} {
        \pgfmathsetmacro{\xx}{(\j-1)*\dx}
        \ifnum\j<3
            \def\fillcol{mint}
        \else
            \def\fillcol{rainbowblue}
        \fi
        \draw[draw=stroke, line width=1pt, fill=\fillcol] (\xx, \yy) circle (\r);

        \pgfmathsetmacro{\drawpurple}{ ((\i+7==0) && (\j==4)) || ( (\i+8==0) && (\j==2))}
        % \pgfmathsetmacro{\drawpurple}{ (\j==2) || (\j==4)}
        \ifnum\drawpurple=1
        \pgfmathsetmacro{\xx}{(\j-1)*\dx}
        \def\fillcol{blue!30!white}
        \draw[draw=stroke, line width=1pt, fill=\fillcol] (\xx, \yy) circle (\r);
        \fi


        \ifnum\i=-7
            \node[anchor=center, font=\large\selectfont] at (\xx, \yy) {$c_{3\j}$};
        \else
            \node[anchor=center, font=\large\selectfont] at (\xx, \yy) {$c_{5\j}$};
        \fi
    %%%%%%%%%%%%%%
    }

    %%%%%%%%%%% Draw left circles
       \foreach \j in {-4.8} {
        \pgfmathsetmacro{\xx}{(\j-1)*\dx}
        \draw[draw=stroke, line width=1pt, fill=green] (\xx, \yy) circle (\r);


        \ifnum\i=-7
            \node[anchor=center, font=\large\selectfont] at (\xx, \yy) {$u_{3}$};
        \else
            \node[anchor=center, font=\large\selectfont] at (\xx, \yy) {$u_{5}$};
        \fi
    %%%%%%%%%%%%%%
    }
}
\end{scope}

\begin{scope}[yshift=-5pt]
% ============
% ============ Draw the bottom basefold code
\foreach \i in {-10} {
    % y coordinate for row
  \pgfmathsetmacro{\yy}{(\i-1.5)*\dy}
  % Circles
    \draw[draw=dashblue, line width=1.1pt, dash pattern=on 5pt off 2pt on 1pt off 2pt, fill=rainbowgray,
    rounded corners=8pt]
    (-6.95, \yy+0.38) rectangle (7.8, \yy-0.38);

    % \ifnum\i=-9
    % \node[anchor=west, font=\large\selectfont] at (-2.4, \yy) {$\Root_3' \Leftarrow \bm{c}_3^{'\,\trans} $};
    % \else
    % \node[anchor=west, font=\large\selectfont] at (-2.4, \yy) {$\Root_5' \Leftarrow \bm{c}_5^{'\,\trans} $};
    % \fi

     % \node[anchor=east, font=\large\selectfont] at (8.0, \yy + 0.8) {$\Enc'(\bm{C}_\mathcal{I})$};

  \foreach \j in {-6,...,9} {
        \pgfmathsetmacro{\xx}{(\j-1)*\dx}
            
        % Decide fill color
        % First 3 columns: top 3 rows cream, bottom 4 rows mint.
        % Last 4 columns: all sky.
          \ifnum\j<-2
            \def\fillcol{mint!90!green}
          \else
            \def\fillcol{rainbowsky}
          \fi

        \draw[draw=stroke, line width=1pt, fill=\fillcol] (\xx, \yy) circle (\r);
        
        % \ifnum\i=-9
        % \node[anchor=center, font=\large\selectfont] at (\xx, \yy) {$c_{3\j}'$};
        % \else
        % \node[anchor=center, font=\large\selectfont] at (\xx, \yy) {$c_{5\j}'$};
        % \fi
        \pgfmathtruncatemacro{\index}{\j+16}
        \node[anchor=center, font=\large\selectfont] at (\xx, \yy) {$c_{\index}'$};
    }
}

%% Matrix U border
\draw[
draw=dashblue,
line width=1.1pt,
dash pattern=on 5pt off 2pt on 1pt off 2pt,
rounded corners=8pt
]
(-2.5,-8.9) rectangle (7.5,-6.8);

\draw[
draw=bordergreen,
line width=1.1pt,
dash pattern=on 5pt off 2pt on 1pt off 2pt,
rounded corners=8pt
]
(-5.7,-8.9) rectangle (-4.95,-6.8);
\end{scope}


% Bottom text: Enc^{\otimes \epsilon^{-1}} \ldots Enc^{\otimes \epsilon^{-1}}
\node[anchor=center, font=\large\selectfont] at (2.3, 1.5) {$\Root\Leftarrow [\Root_1,\ldots,\Root_6]$};
% \node[anchor=center, font=\large\selectfont] at (-6.5, 1.5) {$\Root'_v \Leftarrow \bm{c}' $};
\node[anchor=center, font=\large\selectfont] at (2.4, -6.45) {$\Enc~~\Enc~~\Enc~~\Enc~~\Enc~~\Enc$};
\node[anchor=center, font=\large\selectfont] at (-4.9, -6.85) {$\bm{u} = \Enc(\bm{v})$};
\node[anchor=center, font=\large\selectfont] at (-5.1, -6.35) {$v = \langle \bm{v}, \bm{l} \rangle$};
% \node[anchor=center, font=\large\selectfont] at (-6.4, 0.75) {$\bm{c}' = \Enc'(\bm{v})$};
\node[anchor=center, font=\large\selectfont] at (0.6, 0.9) {$t$};
\node[anchor=center, font=\large\selectfont] at (6.6, -0.4) {$k$};
\node[anchor=center, font=\large\selectfont] at (-7.5, -0.4) {$\bm{v}$};
\node[anchor=center, font=\large\selectfont] at (-2.65, -2.1) {$\mu$};
\node[anchor=center, font=\large\selectfont] at (-4.8, -9.7) {$\Root' \Leftarrow \bm{c}' = \Enc'(\bm{v}||\bm{c}_3^\trans||\bm{c}_5^\trans)$};
\node[anchor=center, font=\large\selectfont] at (6.8, -7.7) {$\bm{C}_\mathcal{I}$};

\node[anchor=center, font=\large\bfseries] at (1.5, -9.6) {CodeSwitch};

% \node[anchor=center, font=\large\selectfont] at (5.7, 1.4) {$v = \bm{l}^{\,\trans}\bm{M}\tilde{\bm{r}}$};
% \node[anchor=center, font=\large\selectfont] at (5.7, 0.6) {$\bm{C} = \bm{G}\bm{M}$};
% \node[anchor=center, font=\large\selectfont] at (5.25, -0.2) {$\Root = \Merkle(\bm{C})$};
% \node[anchor=center, font=\large\selectfont] at (6.5, -0.8) {$= \Merkle([\Root_1, \ldots, \Root_7])$};

\node[anchor=center, font=\Large\selectfont\bfseries, color=titleblue] at (-1.5, 2.0) {TensorSwitch+};
\node[anchor=center, font=\Large\selectfont\bfseries, color=titleblue] at (7.2, 3.0) {DualSwitch};

\draw[-Stealth] (-1,0.7) -- (-4,0.7)
  node[midway, above, font=\Large] {$\bm{v} = \bm{M}\tilde{\bm{r}}$};
\draw[-Stealth] (-1,-6.6) -- (-4,-6.6)
  node[midway, above, font=\Large] {$\bm{u} = \bm{C}\tilde{\bm{r}}$};
  \draw[-Stealth] (-2.7,-8.1) -- (-4.7,-8.1)
  node[midway, above, font=\large] {$\bm{u}_\mathcal{I} = \bm{C}_\mathcal{I}\,\tilde{\bm{r}}$};
\draw[-Stealth, dashed] (-6.4,-8.9) -- (-5.8,-9.4);
\draw[-Stealth, dashed] (-0.5,-9.9) -- (-2.3,-9.7);


\draw[
  decorate,
  decoration={brace, amplitude=6pt}
] (-0.3,0.5) -- (1.3,0.5);
\draw[
  decorate,
  decoration={brace, amplitude=6pt}
] (6.2,0.4) -- (6.2,-2.1);
\draw[
  decorate,
  decoration={brace, amplitude=6pt}
] (-2.2,-4.9) -- (-2.2,0.4);
\draw[
  decorate,
  decoration={brace, amplitude=6pt}
] (-7.0,-2.1) -- (-7.0,0.4);
% \draw[
%   decorate,
%   decoration={brace, amplitude=6pt}
% ] (-2.5,-10.7) -- (-2.5,-9.2);
\end{scope}



\end{tikzpicture}
}
\caption{
Graphical overview of \methodname that integrates \blazeplus with \tensorswitchplus.
$\Enc$ and $\Enc'$ are encoding functions for linear codes $\Code$ and $\Code'$, respectively.
$\Code$ admits efficient encoding, whereas $\Code'$ supports efficient MLE queries.
In \blazeplus, a \emph{structured} vector $\tilde{\bm{r}}$ is used to batch an interleaved codeword $\bm{C}$ into a single codeword $\bm{u}=\bm{C}\bm{r}$, which encodes $\bm v=\bm M\tilde{\bm{r}}$.
Blaze proves proximity of $\bm u$ and simultaneously verifies the value
$v=\langle\bm v,\bm l\rangle$ by switching to the codeword $\bm c'=\Enc'(\bm v)$.
\methodname instantiates the second encoder $\Enc'$ in \blazeplus as the encoder of a tensor code,
whose MLE queries are supported by \tensorswitchplus,
introducing new encoders $\Enc$ and $\Enc'$.
In \tensorswitchplus, the codeword $\bm C$ is similarly batched into
$\bm u=\bm C\tilde{\bm{r}}$, and proximity of $\bm u$ is checked by sampling rows indexed by a small set $\mathcal I$.
Restricting $\bm C$ to $\mathcal I$ yields an interleaved codeword $\bm C_{\mathcal I}$, which, together with the codeword of $\bm v$, is switched into $\Code'$.
}
\label{fig:enc}
\vspace{-10pt}
\end{figure}

\subsection{Tensor-Code Switching via Double Interleaved-Code Switching}

\subsection{Code Switching from Interleaved Code to Tensor Code}

